\subsection{Admissible Injective Projections Have Trivial Holonomy}
  \label{subsec:injective-holonomy}

  This appendix supplies the fiber obstruction used in Section~\ref{sec:necessity-noninjective}.
  In the admissible principal-bundle setting, injectivity makes the structure group, holonomy and connection
  curvature trivial and removes vertical re-identification symmetries.
  These conclusions concern the fibers of $\Pi$; they do not exclude effective base geometry or
  diffeomorphism covariance.

  \subsubsection*{Setup: admissible projections as fibered structures}
    \label{subsec:stepA-setup}

    We work in the regime where $\mathcal{C}_\chi$ and $\mathcal{M}$ admit smooth structures
    (the continuum limit of Section~E.1 and Appendix~D.5).
    In this regime, the projection $\Pi : \mathcal{C}_\chi \to \mathcal{M}$ is required to be \emph{admissible}.

    \begin{definition}[Admissible projection]
      \label{def:admissible-projection}
      A smooth surjective map $\Pi : \mathcal{C}_\chi \to \mathcal{M}$ is called admissible if:
      \begin{enumerate}[label=(A\arabic*)]
  \item \label{A1}
  $\Pi$ is a fiber bundle map.
  For each $y \in \mathcal{M}$, the preimage $\Pi^{-1}(y) \subset \mathcal{C}_\chi$ is a smooth submanifold.
  Fibers are mutually diffeomorphic.
  \item \label{A2}
  $\Pi$ is compatible with the relaxation rule $F$.
  If $x, x' \in \mathcal{C}_\chi$ satisfy $\Pi(x) = \Pi(x')$, then $\Pi(F(x)) = \Pi(F(x'))$.
  Effective dynamics is therefore well-defined on $\mathcal{M}$.
  \item \label{A3}
  $\Pi$ supports a connection.
  There exists a smooth horizontal distribution $H \subset T\mathcal{C}_\chi$ complementary to the vertical,
  with the usual equivariance properties whenever a fiber group action is present.
      \end{enumerate}
    \end{definition}

    Condition~\ref{A1} places $\Pi$ in the category of fiber bundles over $\mathcal{M}$.
    In the principal bundle case, where $G_\Pi$ acts freely and transitively on fibers,
    non-injectivity of $\Pi$ corresponds to $G_\Pi$ being non-trivial.
    For general fiber bundles where the fiber is not identified with the structure group,
    non-injectivity is the weaker condition $|\Pi^{-1}(y)| > 1$.
    The structure group may still be non-trivial even when fibers are discrete.
    The arguments below apply to the principal bundle case, which is the relevant one
    for Cosmochrony (Section~\ref{sec:geometry-pi}, $\Pi \cong S^3$ with $G_\Pi \supseteq \mathrm{SU}(2) \times \mathrm{U}(1)$).

  \subsubsection*{Holonomy as the obstruction to horizontal re-identification}
    \label{subsec:stepA-holonomy}

    Let $\Pi$ be an admissible projection in the principal specialization: $\Pi$ is a smooth principal
    $G_\Pi$-bundle over a connected smooth manifold $\mathcal{M}$, and the horizontal distribution of~\ref{A3} is
    the kernel of a smooth principal connection, a $\mathfrak{g}_\Pi$-valued one-form $\omega$ on $\mathcal{C}_\chi$.
    The connection is fixed throughout; holonomy is a property of $\omega$, not of the bundle alone.
    For a piecewise smooth path $\gamma$ in $\mathcal{M}$ and a point $x$ over $\gamma(0)$, the horizontal lift of
    $\gamma$ starting at $x$ exists and is unique; its endpoint defines the parallel transport $P_\gamma(x)$.
    The holonomy group $\mathrm{Hol}(\omega, x_0)$ at $x_0 \in \mathcal{C}_\chi$, with $y_0 = \Pi(x_0)$, is the
    subgroup of $G_\Pi$ of the elements $g_\gamma$ defined by $P_\gamma(x_0) = x_0\, g_\gamma$ for loops $\gamma$
    based at $y_0$:
    \begin{equation}
      \mathrm{Hol}(\omega, x_0)
      =
      \bigl\{ g_\gamma \in G_\Pi \;\big|\;
      \gamma : [0,1] \to \mathcal{M},
      \;
      \gamma(0) = \gamma(1) = y_0
      \bigr\}
      \subseteq G_\Pi.
    \end{equation}
    The restricted holonomy group $\mathrm{Hol}^0(\omega, x_0)$ is the subgroup obtained from null-homotopic loops.

    Holonomy measures the obstruction to a path-independent identification of fibers by parallel transport for the
    chosen connection.
    It detects the failure of horizontal lifts of loops in $\mathcal{M}$ to close in $\mathcal{C}_\chi$.
    Three objects must be kept apart:
    an ordinary global section, which exists exactly when the bundle is trivial;
    a global horizontal section, a section $s$ with $s^*\omega = 0$, which depends on $\omega$;
    and vertical transformations of $\mathcal{C}_\chi$, among them the gauge transformations.

    \begin{definition}[Re-identification symmetry]
      \label{def:re-id-symmetry}
      A \emph{re-identification symmetry} of the projection $\Pi$ is a transformation
      $\phi : \mathcal{C}_\chi \to \mathcal{C}_\chi$ that:
      \begin{enumerate}[label=(\roman*)]
  \item preserves effective states: $\Pi(\phi(x)) = \Pi(x)$ for all $x$,
  \item is non-trivial: $\phi \neq \mathrm{id}$ on a non-empty open set,
  \item is locally implementable: $\phi$ acts fiberwise and continuously.
      \end{enumerate}
    \end{definition}

    Definition~\ref{def:re-id-symmetry} does not require invertibility, smoothness or $G_\Pi$-equivariance.
    In the principal case, the re-identification symmetries that are smooth, invertible and commute with the
    right $G_\Pi$-action are the non-identity vertical bundle automorphisms, that is, the non-trivial gauge
    transformations $\phi(x) = x\, f(x)$; those that also preserve $\omega$, together with the identity, form a
    subgroup.
    The existence of such transformations is not controlled by holonomy (see the remark on trivialization,
    horizontality and gauge transformations below).

    \begin{proposition}[Trivial holonomy and global horizontal sections]
      \label{prop:holonomy-gauge}
      Let $\Pi : \mathcal{C}_\chi \to \mathcal{M}$ be a smooth principal $G_\Pi$-bundle over a connected smooth
      manifold, with smooth principal connection $\omega$, and let $x_0 \in \mathcal{C}_\chi$.
      Then $\mathrm{Hol}(\omega, x_0) = \{e\}$ if and only if there exists a smooth global horizontal section
      $s : \mathcal{M} \to \mathcal{C}_\chi$, $s^*\omega = 0$, with $s(\Pi(x_0)) = x_0$.
      In that case $(y, g) \mapsto s(y)\, g$ is a global trivialization of $\Pi$ in which the local connection form
      $s^*\omega$ vanishes, and the holonomy at every point of $\mathcal{C}_\chi$ is trivial.
      Consequently, any holonomy group different from $\{e\}$, of any dimension, excludes a global horizontal
      section; it does not exclude an ordinary global section.
    \end{proposition}

    \begin{proof}
      Write $y_0 = \Pi(x_0)$.
      Assume $\mathrm{Hol}(\omega, x_0) = \{e\}$.
      Since $\mathcal{M}$ is connected, every $y \in \mathcal{M}$ is joined to $y_0$ by a piecewise smooth path
      $\gamma$; set $s(y) = P_\gamma(x_0)$.
      If $\gamma'$ is another such path, following $\gamma$ and then $\gamma'$ backwards is a loop based at
      $y_0$; since parallel transport is compatible with concatenation and with the right action, trivial
      holonomy gives $P_{\gamma'}(x_0) = P_\gamma(x_0)$, so $s$ is well defined and $s(y_0) = x_0$.
      On a coordinate ball $B$ centred at $c$, fix one path from $y_0$ to $c$ and use the radial segments from
      $c$; the horizontal lift is the solution of an ordinary differential equation with smooth dependence on
      the endpoint, so $s$ is smooth on $B$.
      For a smooth curve $\sigma$ in $\mathcal{M}$, well-definedness allows $s(\sigma(t))$ to be computed along
      a path to $\sigma(0)$ followed by $\sigma|_{[0,t]}$, so $s \circ \sigma$ is a horizontal curve; hence
      $s^*\omega = 0$.
      Conversely, if $s$ is a horizontal section with $s(y_0) = x_0$, then $s \circ \gamma$ is the horizontal lift
      of every path $\gamma$ starting at $y_0$; for a loop it ends at $s(y_0) = x_0$, so $g_\gamma = e$.
      The map $(y, g) \mapsto s(y)\, g$ is a smooth $G_\Pi$-equivariant bijection with smooth inverse
      $x \mapsto (\Pi(x), g(x))$, where $x = s(\Pi(x))\, g(x)$; this is the trivialization.
      Finally, $\mathrm{Hol}(\omega, x_0\, h) = h^{-1}\,\mathrm{Hol}(\omega, x_0)\, h$, and any two points of
      $\mathcal{C}_\chi$ are related by parallel transport and the right action, so trivial holonomy at $x_0$ is
      trivial holonomy everywhere.
      The last statement follows from the first and from the triviality of the holonomy at every point,
      together with example~(a) of the remark below.
    \end{proof}

    The pulled-back form $s^*\omega$ vanishes, but $\omega$ itself does not vanish when $G_\Pi \neq \{e\}$:
    on vertical vectors it restricts to the Maurer\textendash Cartan form of $G_\Pi$.
    Connectedness is needed for a conclusion from a single basepoint.
    On a disconnected base, trivial holonomy at $x_0$ gives a horizontal section over the component of $\Pi(x_0)$
    only; a global statement requires trivial holonomy at one point over each component.

    \begin{remark}[Trivialization, horizontality and gauge transformations]
      Holonomy neither determines the topology of the bundle nor the existence of gauge transformations.
      Take $\mathcal{C}_\chi = \mathbb{R}^2 \times \mathrm{U}(1)$ over $\mathcal{M} = \mathbb{R}^2$, with fibre
      coordinate $z$; conditions~\ref{A1}--\ref{A3} hold with relaxation rule $F = \mathrm{id}$.
      \begin{enumerate}[label=(\alph*)]
        \item \emph{Trivial bundle, non-trivial holonomy.}
        For $\omega = z^{-1} dz + i\,x\,dy$ the curvature is $i\,dx \wedge dy$.
        Transport around a counterclockwise rectangle of area $a$ multiplies $z$ by $e^{-ia}$, so
        $\mathrm{Hol}(\omega, x_0) = \mathrm{U}(1)$.
        There is no global horizontal section, yet $s_0(x,y) = (x,y,1)$ is an ordinary global section.
        \item \emph{Trivial holonomy, non-trivial gauge transformations.}
        For $\omega = z^{-1} dz$ the holonomy is trivial, while $\phi_f(m, z) = (m, f(m)\, z)$ is a non-trivial
        vertical bundle automorphism for every smooth $f : \mathbb{R}^2 \to \mathrm{U}(1)$ with $f \neq 1$ on a
        non-empty open set; a constant $f \neq 1$ also preserves $\omega$.
        The same constant rotations preserve the curved connection of~(a).
        \item \emph{Discrete holonomy.}
        Over the circle, $\mathcal{C}_\chi = S^1 \times \mathrm{U}(1)$ with $\omega = z^{-1} dz + \tfrac{i}{2}\,d\theta$
        is flat, and the generating loop has holonomy $-1$: the horizontal solution $z(\theta) = e^{-i\theta/2}$ is
        not $2\pi$-periodic.
        The holonomy group $\{\pm 1\}$ is zero-dimensional and still excludes a global horizontal section, while
        $z = 1$ is an ordinary global section.
      \end{enumerate}
      Thus non-trivial holonomy is neither necessary for gauge transformations nor sufficient for a topologically
      non-trivial bundle, and positive dimension of the holonomy group is not the criterion for the absence of a
      global horizontal section.
    \end{remark}

    \begin{remark}[Flat connections]
      The Ambrose\textendash Singer theorem~\cite{AmbroseSinger1953} identifies the Lie algebra of the restricted
      holonomy group with the span of the curvature values $\Omega(u,v)$, for horizontal vectors $u, v$ at points
      joined to $x_0$ by horizontal curves.
      A flat connection ($\Omega = 0$) therefore has trivial restricted holonomy, and its holonomy is the image of
      a representation of $\pi_1(\mathcal{M}, y_0)$ into $G_\Pi$; it is trivial when $\mathcal{M}$ is simply
      connected, and it may be non-trivial otherwise (example~(c) of the preceding remark).
      Ambrose\textendash Singer is used here only for this statement; Proposition~\ref{prop:holonomy-gauge} is
      proved directly above.
      Under injectivity, neither result is needed: $G_\Pi = \{e\}$
      (Theorem~\ref{thm:injective-trivial-holonomy}), so the holonomy group is trivial regardless of the topology
      of $\mathcal{M}$.
    \end{remark}

  \subsubsection*{The triviality theorem for injective admissible projections}
    \label{subsec:stepA-triviality}

    \begin{theorem}[Injective admissible projections have trivial holonomy]
      \label{thm:injective-trivial-holonomy}
      Let $\Pi : \mathcal{C}_\chi \to \mathcal{M}$ be an admissible projection in the sense of
      Definition~\ref{def:admissible-projection}.
      In the principal bundle case, where $G_\Pi$ acts freely and transitively on fibers,
      non-injectivity of $\Pi$ corresponds to $G_\Pi$ being non-trivial.
      For general fiber bundles where the fiber is not identified with the structure group,
      non-injectivity is the weaker condition $|\Pi^{-1}(y)| > 1$.
      The structure group may still be non-trivial even when fibers are discrete.
      The arguments below apply to the principal bundle case, which is the relevant one
      for Cosmochrony (Section~\ref{sec:geometry-pi}, $\Pi \cong S^3$ with
      $G_\Pi \supseteq \mathrm{SU}(2) \times \mathrm{U}(1)$).
      If $\Pi$ is injective, then:
      \begin{enumerate}[label=(\alph*)]
  \item \label{thm-a}
  the fiber over every $y \in \mathcal{M}$ is a single point: $\Pi^{-1}(y) = \{x_y\}$,
  \item \label{thm-b}
  in the principal specialization, the structure group is trivial: $G_\Pi = \{e\}$,
  \item \label{thm-c}
  in that principal case, the holonomy group is trivial:
  $\mathrm{Hol}(\omega, x_0) = \{e\}$ for any connection $\omega$ and any basepoint $x_0$,
  \item \label{thm-d}
  no non-trivial re-identification symmetry exists in the sense of Definition~\ref{def:re-id-symmetry}.
      \end{enumerate}
    \end{theorem}

    \begin{proof}
      Part~\ref{thm-a} is immediate from injectivity.
      If $\Pi(x) = \Pi(x') = y$, then $x = x'$, hence $|\Pi^{-1}(y)| = 1$ for all $y$.

      Part~\ref{thm-b} is a principal-bundle statement.
      If fibers are singletons, the only group acting freely and transitively on a one-element set is $\{e\}$.

      Part~\ref{thm-c} follows from $\mathrm{Hol}(\omega, x_0) \subseteq G_\Pi$.
      If $G_\Pi = \{e\}$, then $\mathrm{Hol}(\omega, x_0) = \{e\}$.

      Part~\ref{thm-d} follows directly from injectivity.
      If $\Pi(\phi(x)) = \Pi(x)$ for all $x$, injectivity implies $\phi(x) = x$ for all $x$.
      Hence $\phi = \mathrm{id}$, contradicting non-triviality.
    \end{proof}

  \subsubsection*{From the trivial structure group to absence of fiber gauge structure}
    \label{subsec:stepA-gauge}

    \begin{corollary}[No emergent gauge field from injective projection]
      \label{cor:no-gauge}
      Under the hypotheses of Theorem~\ref{thm:injective-trivial-holonomy} in the principal specialization,
      any connection is trivial in the sense that its curvature vanishes.
      No non-trivial gauge field can emerge from the projection structure alone.
    \end{corollary}

    \begin{proof}
      With $G_\Pi = \{e\}$, the Lie algebra is $\{0\}$.
      Any connection is therefore $\omega = 0$, and $\Omega = d\omega + \omega \wedge \omega = 0$.
      A non-trivial gauge field cannot be obtained as the connection of this trivial principal structure.
      This does not exclude other constructions from $(F,\Pi)$.
    \end{proof}

    \begin{corollary}[No vertical re-identification under injectivity]
      \label{cor:no-diffeo}
      Under the hypotheses of Theorem~\ref{thm:injective-trivial-holonomy}, every fiber-preserving
      transformation $\phi$ with $\Pi\circ\phi=\Pi$ is the identity.
    \end{corollary}

    \begin{proof}
      Injectivity gives $\phi(x)=x$ from $\Pi(\phi(x))=\Pi(x)$ for every $x$.
    \end{proof}

    This vertical statement supplies no obstruction to D2.
    Even for a non-injective projection, vertical transformations fix all base states and therefore do not
    generate base diffeomorphisms.
    Redundancy among metric representatives on the base is a different quotient.

  \subsubsection*{Fiber curvature and base curvature: status of the bridge}
    \label{subsec:stepA-curvature}

    The preceding results exclude nontrivial vertical re-identification under injectivity.
    They do not exclude a curved effective metric on the base.
    This subsection records what is, and what is not, established about a contribution of fiber curvature to the
    curvature of the effective base metric.

    In the Cosmochrony framework, the effective metric on $\mathcal{M}$
    is defined by the Hessian of the effective Lagrangian
    (Appendix~\ref{sec:born-lagrangian_derivation}, Eq.~\eqref{eq:A12-emergent-metric_v171}):
    \begin{equation}
      g^\mathrm{eff}_{\mu\nu}
      \propto
      \frac{\partial^2 \mathcal{L}_\mathrm{eff}}
      {\partial(\partial_\mu \chi)\,\partial(\partial_\nu \chi)},
      \label{eq:eff-metric-hessian}
    \end{equation}
    valid in projectable regimes.
    This expression involves neither the connection $\omega$ nor any metric on the total space $\mathcal{C}_\chi$.
    By itself, it therefore provides no mechanism through which the fiber curvature $\Omega$ enters $g^\mathrm{eff}$.

    A fiber-to-base curvature relation of Kaluza--Klein type exists only for a different object, and only under
    additional hypotheses that are not among the admissibility conditions~\ref{A1}--\ref{A3}:
    \begin{enumerate}[label=(K\arabic*)]
      \item \label{K1}
      $\mathcal{C}_\chi$ carries the metric $\hat g = \Pi^* g + q(\omega, \omega)$, where $g$ is a metric on
      $\mathcal{M}$ and $q$ is an $\mathrm{Ad}$-invariant inner product on $\mathfrak{g}_\Pi$, independent of the
      base point; $\Pi$ is then a Riemannian submersion with totally geodesic fibers, bi-invariant fiber metric,
      and horizontal distribution $\ker\omega$;
      \item \label{K2}
      the effective metric on $\mathcal{M}$ is identified with the base metric of that submersion;
      \item \label{K3}
      for a sourcing reading, a field equation is imposed on the total space (for instance
      $\widehat{\mathrm{Ric}} = 0$).
    \end{enumerate}

    \begin{proposition}[Conditional fiber-to-base curvature relation]
      \label{prop:fiber-base-curvature}
      Let $\Pi : \mathcal{C}_\chi \to \mathcal{M}$ be an admissible principal $G_\Pi$-bundle with connection $\omega$
      and curvature $\Omega$, and assume~\ref{K1} and~\ref{K2}.
      Write $\Omega = \Omega^a T_a$ in a basis $(T_a)$ of $\mathfrak{g}_\Pi$, with $q_{ab} = q(T_a, T_b)$, and use the
      convention in which the Ricci tensor is positive on round spheres.
      Then, on horizontal lifts of vectors tangent to $\mathcal{M}$, the base Ricci tensor and the total-space Ricci
      tensor differ by a term quadratic in the curvature,
      \begin{equation}
        R^\mathrm{eff}_{\mu\nu}
        =
        \widehat{R}^{\,\mathrm{hor}}_{\mu\nu}
        + \frac{1}{2}\, q_{ab}\,\Omega^a_{\mu\alpha}\,\Omega^b{}_{\nu}{}^{\alpha},
        \label{eq:fiber-to-base-curvature}
      \end{equation}
      where the index $\alpha$ is raised with $g$.
      The contraction uses the inner product $q$ and not a trace in a chosen representation; in particular it does not
      vanish for $G_\Pi = \mathrm{U}(1)$.
    \end{proposition}

    \begin{proof}
      Choose a $q$-orthonormal basis $(T_a)$, so that the structure constants $f_{abc}$ are totally antisymmetric
      by $\mathrm{Ad}$-invariance, and use the orthonormal coframe $(e^i, e^a)$ of $\hat g$,
      where $(e^i)$ is the pull-back of an orthonormal coframe of $g$ with Levi-Civita forms $\omega^i{}_j$,
      and $e^a = \omega^a$.
      The structure equation $d\omega = \Omega - \tfrac12[\omega,\omega]$ gives
      $de^a = \tfrac12\Omega^a_{ij}\,e^i\wedge e^j - \tfrac12 f^a{}_{bc}\,e^b\wedge e^c$
      and $de^i = -\omega^i{}_j\wedge e^j$.
      The torsion-free, metric connection forms of $\hat g$ are therefore
      \begin{equation*}
        \hat\omega^i{}_j = \omega^i{}_j - \tfrac12\,\Omega^a_{ij}\,e^a,
        \qquad
        \hat\omega^a{}_i = -\hat\omega^i{}_a = \tfrac12\,\Omega^a_{ij}\,e^j,
        \qquad
        \hat\omega^a{}_b = -\tfrac12\, f^a{}_{bc}\,e^c .
      \end{equation*}
      Computing the horizontal components of the Ricci tensor from these forms, the terms involving $f_{abc}$ drop out,
      because they enter only through $f^a{}_{ba} = 0$ or through vertical--vertical components.
      What remains is the abelian Kaluza--Klein computation, applied component by component, which gives
      $\widehat{R}_{ij} = R_{ij} - \tfrac12\,\Omega^a_{ik}\,\Omega^a_{jk}$ in the orthonormal frame, that is,
      Eq.~\eqref{eq:fiber-to-base-curvature}.
      For $G_\Pi = \mathrm{U}(1)$ with $\hat g = g + \phi^2(dz + A)^2$ over a two-dimensional base with metric
      $e^{2f}(dx^2 + dy^2)$, the coefficient $\tfrac12$ was checked by direct symbolic computation of the Christoffel
      symbols of $\hat g$ (script \texttt{code/kk\_horizontal\_ricci\_check.py} in the paper repository).
      Hypothesis~\ref{K2} is what allows the base metric to be read as the effective metric.
    \end{proof}

    \begin{remark}[Status]
      Equation~\eqref{eq:fiber-to-base-curvature} is an identity between two Ricci tensors.
      It generates base curvature from fiber curvature only together with a field equation on the total space,
      as in~\ref{K3}.
      Hypotheses~\ref{K1}--\ref{K3} are not derived in this paper, and~\ref{K2} is a different definition of the
      effective metric from Eq.~\eqref{eq:eff-metric-hessian}; their mutual consistency is not shown.
      Earlier versions of this appendix presented the relation as following from Eq.~\eqref{eq:eff-metric-hessian}
      through an unspecified ``induction scheme''.
      That argument is withdrawn.
      The missing element is a construction of $\hat g$, or of another explicit relation, from $(F, \Pi)$; no result
      here shows that such a construction cannot exist.
    \end{remark}

    \begin{remark}[The injective case]
      In the principal specialization,
      \begin{equation}
        \Pi \text{ injective}
        \Rightarrow
        G_\Pi = \{e\}
        \Rightarrow
        \Omega = 0
        \Rightarrow
        \text{no fiber-induced curvature term on } \mathcal{M}.
        \label{eq:injective-chain}
      \end{equation}
      The implication runs from the trivial group to the vanishing curvature (Corollary~\ref{cor:no-gauge}); a flat
      connection does not by itself make the structure group trivial.
      With singleton fibers there is no fiber geometry, so every fiber-induced term vanishes by definition.
      Equation~\eqref{eq:injective-chain} therefore adds nothing to Theorem~\ref{thm:injective-trivial-holonomy}
      and Corollary~\ref{cor:no-gauge}, and it does not depend on Proposition~\ref{prop:fiber-base-curvature}.
      It holds in the principal case; the general fiber-bundle case is not covered.
    \end{remark}

  \subsubsection*{The minimal non-injective case: $G_\Pi = \mathrm{U}(1)$}
    \label{subsec:stepA-minimal}

    As a consistency check and to connect to Section~\ref{sec:geometry-pi},
    we consider the minimal non-trivial case $G_\Pi = \mathrm{U}(1)$.
    A principal $\mathrm{U}(1)$-bundle over $\mathcal{M}$ with connection $\omega$ has:
    \begin{align}
      \mathrm{Hol}(\omega, x_0) &\subseteq \mathrm{U}(1), \\
      \Omega &= d\omega \in \Omega^2(\mathcal{M}, i\mathbb{R}),
    \end{align}
    where $\Omega$ is interpreted as the electromagnetic field strength in the Cosmochrony setting (Section~A.11).
    The first Chern class $c_1(\Pi) = [\Omega / 2\pi i] \in H^2(\mathcal{M}, \mathbb{Z})$
    is non-trivial if and only if the
    bundle is topologically non-trivial.

    This matches the structure of Section~\ref{sec:gauges-transmittance}.
    The electromagnetic $\mathrm{U}(1)$ gauge symmetry emerges because the projection fiber supports a non-trivial
    Hopf fibration $S^1 \hookrightarrow S^3 \to S^2$.

    In the injective limit, where the fiber degenerates to a point, the bundle trivializes and the $\mathrm{U}(1)$ gauge
    structure disappears, consistent with Eq.~\eqref{eq:injective-chain}.

  \subsubsection*{Summary of Step A}
    \label{subsec:stepA-summary}

    The following table summarizes the principal-bundle fiber obstruction used in
    Proposition~\ref{thm:necessity} and the additional geometric constructions discussed above.

    \begin{center}
      \begin{tabular}{p{0.26\linewidth}p{0.24\linewidth}p{0.38\linewidth}}
        \hline
        \raggedright Fiber type & \raggedright Structure group $G_\Pi$ & \raggedright Emergent structure
        \tabularnewline
        \hline
        \raggedright Point ($|\Pi^{-1}(y)| = 1$, injective)
        & \raggedright $G_\Pi = \{e\}$
        & \raggedright No fiber gauge field, no fiber-induced curvature term, no vertical re-identification
        \tabularnewline
        \raggedright Circle ($S^1$, minimal non-injective)
        & \raggedright $G_\Pi = \mathrm{U}(1)$
        & \raggedright $\mathrm{U}(1)$ gauge field, electromagnetic curvature
        \tabularnewline
        \raggedright Three-sphere ($S^3$, Cosmochrony setting)
        & \raggedright $G_\Pi \supseteq \mathrm{SU}(2) \times \mathrm{U}(1)$
        & \raggedright Gauge structure; base-curvature term only under~\ref{K1}--\ref{K3}
        (Prop.~\ref{prop:fiber-base-curvature})
        \tabularnewline
        \hline
      \end{tabular}
    \end{center}

    In the principal case, injectivity removes vertical re-identification symmetry and fiber-induced curvature.
    Conversely, non-injectivity alone is not sufficient for nontrivial holonomy, nonzero connection curvature,
    or a gravitational response.
    Whether a non-injective projection sources effective base curvature requires hypotheses~\ref{K1}--\ref{K3}
    or another explicit construction (Proposition~\ref{prop:fiber-base-curvature}).
    A base metric defined independently of fiber curvature, including the Hessian proposal above, is not
    ruled out under injectivity.
    Appendix~\ref{subsec:IR} records a separate open question: construction of the infrared gravitational action.
